\section{Algorithm}
\label{sec:algorithm}

This section specifies the agentic system built on \textsc{DeXposure-FM} for two tasks: monitoring and decision-making. 

\subsection{Task I: Network Risk Monitoring}
\label{subsec:task_monitoring}

At each decision epoch $t$, we observe a current DeFi credit-exposure snapshot $G_t=(V,E_t,W_t)$ and node covariates $X_t$ (e.g., log-TVL and auxiliary features). Given a set of forecast horizons $\mathcal{H}$, the foundation model \textsc{DeXposure-FM} induces predictive distributions over future exposure networks,
\begin{equation}
\mathcal{P}_{t,h} \;:=\; p\!\left(G_{t+h}\mid G_{\le t}, X_{\le t}\right), \qquad h\in\mathcal{H}.
\end{equation}
We denote by $\widehat{G}_{t+h}$ a representative predicted network constructed from $\mathcal{P}_{t,h}$ (e.g., the expected-weight graph), and by $\widetilde{G}^{(m)}_{t+h}\sim\mathcal{P}_{t,h}$ a sampled predicted network for uncertainty quantification.

\paragraph{Monitoring objective.}
The monitoring agent computes a collection of systemic indicators (``measurements'') as functionals of predicted networks:
\begin{equation}
\mathbf{M}_{t,h} \;:=\; \Phi\!\left(\widehat{G}_{t+h}, X_t\right)\in\mathbb{R}^d,
\qquad
\widetilde{\mathbf{M}}^{(m)}_{t,h} \;:=\; \Phi\!\left(\widetilde{G}^{(m)}_{t+h}, X_t\right),
\end{equation}
where $\Phi$ may include, for example, concentration indices, sector spillover measures, systemic importance scores, and stress-test losses. Let $\mathcal{B}_t(\cdot)$ be a baseline band operator derived from historical measurements (e.g., rolling mean and standard deviation). The monitoring agent issues an alert when a predicted measurement departs from its baseline beyond a tolerance:
\begin{equation}
\label{eq:monitor_alert_rule}
\mathcal{A}_t \;:=\; \Bigl\{\, a=(h,j): \; g\!\left(M_{t,h}^{(j)}, \mathcal{B}_t^{(j)}\right)=1 \Bigr\},
\end{equation}
where $M_{t,h}^{(j)}$ is the $j$-th component of $\mathbf{M}_{t,h}$, $\mathcal{B}_t^{(j)}$ is the corresponding baseline band, and $g(\cdot)$ is a trigger rule (e.g., a $z$-score or change-point test). 

\paragraph{Uncertainty-aware confidence.}
To expose forecast uncertainty, the agent attaches an alert confidence score. Let $\mathrm{U}_{t,h}^{(j)}$ denote a dispersion statistic (e.g., inter-quantile range) computed from $\{\widetilde{\mathbf{M}}^{(m)}_{t,h}\}_{m=1}^M$, and let $DH_t\in[0,1]$ be a data-health score. We define a generic confidence mapping
\begin{equation}
\label{eq:monitor_confidence}
C_t(a)\;=\;\psi\!\left(DH_t,\, \mathrm{U}_{t,h}^{(j)},\, h\right)\in[0,1],
\end{equation}
so alerts carry both evidence (triggering statistics) and reliability (data/model quality).


\subsection{Task II: Decision-Making}
\label{subsec:task_decision}

The decision-making agent consumes monitoring alerts $\mathcal{A}_t$ and predictive distributions $\{\mathcal{P}_{t,h}\}_{h\in\mathcal{H}}$ to propose risk-mitigating actions. Let $\mathcal{U}$ denote a feasible action set (e.g., watchlist updates, exposure reduction recommendations, governance parameter suggestions). An action $u\in\mathcal{U}$ induces an intervention operator on the current state,
\begin{equation}
T(u): (G_t, X_t)\mapsto (G_t^{u}, X_t^{u}),
\end{equation}
representing, for example, reduced exposure to certain protocols, altered risk limits, or scenario-dependent constraints. Under $u$, the foundation model yields counterfactual predictive distributions
\begin{equation}
\mathcal{P}^{u}_{t,h} \;:=\; p\!\left(G_{t+h}\mid G_{\le t}^{u}, X_{\le t}^{u}\right), \qquad h\in\mathcal{H}.
\end{equation}
Let $\mathcal{S}$ be a library of standardised stress scenarios (shocks) applied to predicted networks, and let $\ell(\cdot)$ be a nonnegative systemic loss functional (e.g., stress-test loss, tail concentration, spillover loss). For a scenario $s\in\mathcal{S}$, define the (random) loss at horizon $h$:
\begin{equation}
L_{t,h}(u,s) \;:=\; \ell\!\left(G_{t+h}^{u,s}\right), \qquad
G_{t+h}^{u,s} \sim \mathcal{P}^{u}_{t,h}\ \text{with scenario } s,
\end{equation}
and an uncertainty-aware risk aggregation (e.g., expectation plus tail risk):
\begin{equation}
\label{eq:risk_objective}
\mathcal{R}_t(u) \;:=\; \sum_{h\in\mathcal{H}} \omega_h \cdot 
\max_{s\in\mathcal{S}} \Bigl(
\mathbb{E}[L_{t,h}(u,s)] + \lambda\cdot \mathrm{CVaR}_{\alpha}(L_{t,h}(u,s))
\Bigr),
\end{equation}
where $\omega_h$ are horizon weights, $\lambda\ge 0$ controls tail sensitivity, and $\alpha\in(0,1)$ is a tail level.

\paragraph{Decision objective with constraints.}
Let $\mathrm{cost}(u)$ be the operational cost of recommending or implementing $u$ (including conservativeness penalties), and let $\mathcal{C}_t$ encode safety constraints such as data-health and calibration gates. The decision-making task is to select a recommendation
\begin{equation}
\label{eq:decision_problem}
u_t^\star \;\in\; \arg\min_{u\in\mathcal{U}} \ \mathcal{R}_t(u) + \eta\cdot \mathrm{cost}(u)
\quad \text{s.t.}\quad u \in \mathcal{C}_t,
\end{equation}
where $\eta\ge 0$ trades off risk reduction and operational cost. In the minimal viable agent, $\mathcal{U}$ is restricted to a small set of conservative playbooks and $\mathcal{C}_t$ enforces strict gating: when $DH_t$ (or calibration) is low, the feasible set collapses to non-interventional actions (e.g., \texttt{Monitor}/\texttt{Investigate} recommendations).

\paragraph{Decision outputs.}
The agent returns a ranked set of recommendation tickets $\mathcal{D}_t=\{(u,r(u))\}$, each with (i) the proposed action $u$, (ii) an estimated risk reduction $r(u)$ derived from \eqref{eq:risk_objective}--\eqref{eq:decision_problem}, and (iii) an evidence bundle summarizing the triggering alerts, scenario impacts, and uncertainty measures.

\subsection{System Interface}
At each epoch $t$, the agent receives the latest state (graph snapshot and features) and returns: (a) a set of monitoring alerts with evidence and confidence, (b) a scenario stress summary, and (c) a ranked list of recommendation tickets. The agent interacts with \textsc{DeXposure-FM} through a single primitive:
\begin{equation}
\label{eq:agent_forecast_primitive}
\mathcal{P}_{t,h} \leftarrow \textsc{DeXposure-FM}.\mathrm{Forecast}(G_t,X_t,h), \forall h\in\mathcal{H},
\end{equation}
where $\mathcal{P}_{t,h}$ denotes the model's predictive distribution for horizon $h$. We use a deterministic graph-construction operator $\mathrm{BuildPredGraph}(\cdot)$ to obtain a representative forecasted network $\widehat{G}_{t+h}$ from $\mathcal{P}_{t,h}$ (e.g., expected-weight or thresholded-probability construction). Optionally, we draw Monte Carlo samples $\{\widetilde{G}^{(m)}_{t+h}\}_{m=1}^M$ for uncertainty quantification.

\subsection{Data-Health Gating}
\label{subsec:data_health_gate}
Before producing actionable outputs, the agent computes a scalar data-health score
\begin{align}
\label{eq:data_health}
&DH_t \leftarrow \mathrm{DataHealth}(G_t,X_t)\in[0,1],\nonumber\\
&\mathrm{SAFE\_MODE}_t := \mathbb{I}\{DH_t<\tau_{\mathrm{data}}\}.
\end{align}
$\mathrm{DataHealth}(\cdot)$ aggregates deterministic checks including freshness, missingness, discontinuities/outliers, and topology sanity. In safe mode, the agent still emits monitoring alerts but suppresses intervention-type recommendations (Section~\ref{subsec:decision_v1}).

\subsection{Monitoring: Forecast--then--Measure--then--Alert}
\label{subsec:monitoring_agent}
For each horizon $h\in\mathcal{H}$, the agent forecasts $\mathcal{P}_{t,h}$ by \eqref{eq:agent_forecast_primitive}, constructs $\widehat{G}_{t+h}$, and evaluates a vector of monitoring functionals:
\begin{equation}
\label{eq:monitor_metrics}
\mathbf{M}_{t,h} := \Phi(\widehat{G}_{t+h},X_t)\in\mathbb{R}^d,~
\widetilde{\mathbf{M}}^{(m)}_{t,h} := \Phi(\widetilde{G}^{(m)}_{t+h},X_t),
\end{equation}
where $\Phi(\cdot)$ includes the systemic indicators used by the monitoring dashboard (e.g., systemic-importance ranking, cross-sector spillover concentration, stability/concentration indices, and stress-test losses). Let $\mathcal{B}_t^{(j)}$ be a baseline band for the $j$-th metric (e.g., a rolling mean/standard deviation band over a fixed historical window). An alert is triggered when a forecasted metric departs from its baseline:
\begin{equation}
\label{eq:alert_trigger_generic}
\mathcal{A}_t :=
\Bigl\{\, a=(h,j)\;:\; g\!\bigl(M_{t,h}^{(j)},\,\mathcal{B}_t^{(j)}\bigr)=1 \Bigr\},
\end{equation}
where $g(\cdot)$ is a deterministic trigger (e.g., $z$-score, change-point score, or rule-based condition).

\paragraph{Evidence and attribution.}
Each alert $a=(h,j)\in\mathcal{A}_t$ is returned with an evidence bundle:
\begin{equation}
\label{eq:alert_evidence}
\mathrm{Ev}_t(a) := \Bigl( M_{t,h}^{(j)},\ \mathcal{B}_t^{(j)},\ \mathrm{Attr}_t(a),\ h \Bigr),
\end{equation}
where $\mathrm{Attr}_t(a)$ denotes a ranked set of contributing nodes/edges/sectors obtained via marginal contribution approximations on the corresponding functional $\Phi^{(j)}(\cdot)$.

\paragraph{Uncertainty-aware confidence.}
Using sampled metric realizations $\{\widetilde{\mathbf{M}}^{(m)}_{t,h}\}_{m=1}^M$, we compute a dispersion statistic $\mathrm{U}_{t,h}^{(j)}$ (e.g., inter-quantile range) and attach a confidence score:
\begin{equation}
\label{eq:alert_confidence_generic}
C_t(a) := \psi\!\left(DH_t,\ \mathrm{U}_{t,h}^{(j)},\ h\right)\in[0,1],
\end{equation}
which monotonically decreases with lower data health, wider predictive spread, and longer horizons.

\subsection{Scenario Engine}
\label{subsec:scenario_engine}
To support systematic triage, the agent evaluates a small standardised scenario set $\mathcal{S}_0$ (stress shocks). For each horizon $h$ and scenario $s$, we compute a scenario loss functional:
\begin{equation}
\label{eq:scenario_loss}
L_{t,h}(s) := \ell\!\left(\widehat{G}_{t+h}; s\right),~
\widetilde{L}^{(m)}_{t,h}(s) := \ell\!\left(\widetilde{G}^{(m)}_{t+h}; s\right),
\end{equation}
where $\ell(\cdot;\,s)$ denotes the stress-test procedure applied to the predicted network under shock $s$. The scenario summary $\mathcal{R}_t$ ranks scenarios by expected (or median) loss and highlights worst-case horizons and dominant propagation channels via the same attribution interface used in monitoring.

\subsection{Decision Agent v1: Conservative Playbook Recommendations}
\label{subsec:decision_v1}
The decision module maps $(\mathcal{A}_t,\{\mathrm{Ev}_t(a),C_t(a)\},\mathcal{R}_t)$ to a ranked set of recommendation tickets. Let $\mathcal{U}$ be a set of conservative recommendation types (e.g., watchlist update, escalated monitoring, investigation request, exposure reduction suggestion, contingency planning). Each $u\in\mathcal{U}$ must satisfy a safety constraint set $\mathcal{C}_t$:
\begin{align}
\label{eq:decision_constraints}
&u \in \mathcal{C}_t
\quad \Longleftrightarrow \nonumber\\
&\Bigl(\mathrm{SAFE\_MODE}_t=0\Bigr)\ \wedge\ \Bigl(\min_{a\in \mathrm{Trig}(u)} C_t(a) \ge \tau_{\mathrm{conf}}\Bigr),
\end{align}
where $\mathrm{Trig}(u)\subseteq\mathcal{A}_t$ denotes the subset of alerts required to justify $u$. In safe mode, the feasible set collapses to non-interventional actions (e.g., \texttt{Monitor}/\texttt{Investigate}).

\paragraph{Ticket ranking.}
Each feasible recommendation is scored by a monotone heuristic combining severity, confidence, and scenario impact:
\begin{equation}
\label{eq:ticket_score}
\mathrm{Score}_t(u) :=
\mathrm{Sev}_t(u)\cdot \mathrm{Conf}_t(u)\cdot \mathrm{Imp}_t(u),
\end{equation}
where $\mathrm{Sev}_t(u)$ is derived from triggered metric deviations, $\mathrm{Conf}_t(u)$ aggregates $C_t(a)$ over required alerts, and $\mathrm{Imp}_t(u)$ is derived from $\mathcal{R}_t$ (e.g., worst-case scenario loss). The decision agent returns a ranked ticket list $\mathcal{D}_t=\{(u,\mathrm{Score}_t(u),\mathrm{Ev}_t(u))\}$, where $\mathrm{Ev}_t(u)$ bundles alert evidence, scenario summaries, and an auditable rationale.

\subsection{End-to-End Procedure}
Algorithm~\ref{alg:agent_loop} summarises the minimal viable agent loop.

\begin{algorithm}[t]
\caption{DeXposure-Agent}
\label{alg:agent_loop}
\begin{algorithmic}[1]
\Require Current state $(G_t,X_t)$; horizons $\mathcal{H}$; scenarios $\mathcal{S}_0$; baseline bands $\{\mathcal{B}_t^{(j)}\}_{j=1}^d$; thresholds $(\tau_{\mathrm{data}},\tau_{\mathrm{conf}})$
\Ensure Alerts $\mathcal{A}_t$ with evidence/confidence; scenario summary $\mathcal{R}_t$; recommendation tickets $\mathcal{D}_t$
\State $DH_t \gets \mathrm{DataHealth}(G_t,X_t)$; $\mathrm{SAFE\_MODE}_t \gets \mathbb{I}\{DH_t<\tau_{\mathrm{data}}\}$
\ForAll{$h\in\mathcal{H}$}
    \State $\mathcal{P}_{t,h} \gets \textsc{DeXposure-FM}.\mathrm{Forecast}(G_t,X_t,h)$
    \State $\widehat{G}_{t+h} \gets \mathrm{BuildPredGraph}(\mathcal{P}_{t,h})$
    \State $\mathbf{M}_{t,h} \gets \Phi(\widehat{G}_{t+h},X_t)$
    \State $\mathcal{A}_{t,h} \gets \{(h,j): g(M^{(j)}_{t,h},\mathcal{B}_t^{(j)})=1\}$
    \State attach $\mathrm{Ev}_t(a)$ and $C_t(a)$ for all $a\in\mathcal{A}_{t,h}$
    \State compute scenario losses $\{L_{t,h}(s)\}_{s\in\mathcal{S}_0}$ and attributions
\EndFor
\State $\mathcal{A}_t \gets \bigcup_{h\in\mathcal{H}}\mathcal{A}_{t,h}$
\State $\mathcal{R}_t \gets \mathrm{SummarizeScenarios}(\{L_{t,h}(s)\}_{h,s})$
\State $\mathcal{D}_t \gets \mathrm{PlaybookDecision}(\mathcal{A}_t,\mathcal{R}_t, DH_t,\mathrm{SAFE\_MODE}_t,\tau_{\mathrm{conf}})$
\State \Return $(\mathcal{A}_t,\mathcal{R}_t,\mathcal{D}_t)$
\end{algorithmic}
\end{algorithm}